\documentclass[11pt]{article}

\usepackage[a4paper,margin=1in]{geometry}
\usepackage{amsmath,amssymb,booktabs}
\usepackage[hidelinks]{hyperref}

\title{Phase 1 Static FJSP Mathematical Model}
\author{FJSP-AGV-RL Research Platform}
\date{\today}

\begin{document}
\maketitle

\section{Scope}

This document describes the Phase 1 static flexible job shop scheduling
problem (FJSP) model implemented by the repository. Transportation time,
AGV routing, dynamic job arrivals, reinforcement learning, and graph neural
network components are outside the Phase 1 scope.

\section{Sets and Indices}

\begin{table}[h]
\centering
\begin{tabular}{ll}
\toprule
Symbol & Meaning \\
\midrule
$J$ & set of jobs, indexed by $j$ \\
$O_j$ & ordered set of operations for job $j$, indexed by $o$ \\
$M$ & set of machines, indexed by $m$ \\
$A_{jo}$ & eligible machines for operation $(j,o)$ \\
\bottomrule
\end{tabular}
\end{table}

Operation $(j,o)$ must precede $(j,o+1)$ for every job $j$ and operation
position $o$ where $o+1 \in O_j$.

\section{Parameters}

\begin{table}[h]
\centering
\begin{tabular}{ll}
\toprule
Symbol & Meaning \\
\midrule
$p_{jom}$ & processing time of operation $(j,o)$ on machine $m$ \\
$H$ & conservative scheduling horizon \\
\bottomrule
\end{tabular}
\end{table}

The implementation uses a conservative horizon
\[
H = \sum_{j \in J} \sum_{o \in O_j} \max_{m \in A_{jo}} p_{jom}.
\]
This upper bound is not necessarily tight, but it is safe for the small
CP-SAT baseline model.

\section{Decision Variables}

\begin{table}[h]
\centering
\begin{tabular}{ll}
\toprule
Symbol & Meaning \\
\midrule
$x_{jom} \in \{0,1\}$ & 1 if operation $(j,o)$ is assigned to machine $m$ \\
$S_{jo} \ge 0$ & start time of operation $(j,o)$ \\
$C_{jo} \ge 0$ & completion time of operation $(j,o)$ \\
$C_{\max} \ge 0$ & makespan \\
$y_{abm} \in \{0,1\}$ & order variable for two operations $a,b$ on machine $m$ \\
\bottomrule
\end{tabular}
\end{table}

Here $a$ and $b$ denote operation identifiers such as $(j,o)$. The explicit
$y$ variables are useful for a classical mixed-integer programming statement.
The current OR-Tools implementation uses optional interval variables and
\texttt{NoOverlap} instead of manually creating all $y$ variables.

\section{Objective}

The objective is to minimize the makespan:
\[
\min C_{\max}.
\]

\section{Constraints}

\subsection{Machine Assignment}

Each operation must be assigned to exactly one eligible machine:
\[
\sum_{m \in A_{jo}} x_{jom} = 1,
\qquad \forall j \in J,\; o \in O_j.
\]

Machines outside $A_{jo}$ are not valid choices and are never created as
candidate assignments in the implementation.

\subsection{Processing Time Link}

If operation $(j,o)$ is assigned to machine $m$, then its completion time must
match its start time plus the selected processing time:
\[
C_{jo} \ge S_{jo} + p_{jom} - H(1 - x_{jom}),
\qquad \forall j,o,\; m \in A_{jo}.
\]

\subsection{Job Precedence}

Operations within the same job must follow the given route order:
\[
C_{jo} \le S_{j,o+1},
\qquad \forall j \in J,\; o,o+1 \in O_j.
\]

\subsection{Machine Capacity}

A machine can process at most one operation at a time. For any two distinct
operations $a=(j,o)$ and $b=(k,q)$ that are both eligible for machine $m$,
the disjunctive ordering is:
\[
S_a + p_{am} \le S_b + H(1 - y_{abm}) + H(2 - x_{am} - x_{bm}),
\]
\[
S_b + p_{bm} \le S_a + Hy_{abm} + H(2 - x_{am} - x_{bm}).
\]

These constraints are active only when both operations are assigned to machine
$m$. In the CP-SAT implementation, this logic is represented by optional
intervals and a per-machine \texttt{NoOverlap} constraint.

\subsection{Makespan Definition}

The makespan must be no smaller than every final operation completion time:
\[
C_{\max} \ge C_{j,last(j)}, \qquad \forall j \in J.
\]

\section{Feasibility Checker Mapping}

The repository's feasibility checker validates the following model conditions:

\begin{itemize}
  \item every operation appears exactly once;
  \item all job and operation ids are inside the instance range;
  \item every selected machine is eligible for the operation;
  \item recorded processing time matches the instance data;
  \item job precedence constraints are satisfied;
  \item no two operations overlap on the same machine;
  \item stored makespan equals the maximum recorded completion time.
\end{itemize}

\section{Dispatching Rules}

The strict dispatching rules use a common candidate tuple
\[
(j,o,m,r_{jo},s_{jom},f_{jom},p_{jom}),
\]
where $r_{jo}$ is the job ready time, $s_{jom}$ is the earliest feasible start
time on machine $m$, and $f_{jom}=s_{jom}+p_{jom}$ is the tentative finish time.

\begin{table}[h]
\centering
\begin{tabular}{lll}
\toprule
Rule & Operation selection & Machine selection \\
\midrule
FIFO & smallest $(r_{jo},j,o)$ & earliest finish \\
SPT & smallest $p_{jom}$ pair & selected with pair \\
EFT & smallest $f_{jom}$ pair & selected with pair \\
MWKR & largest remaining route work & earliest finish \\
Random & random ready pair & selected with pair \\
\bottomrule
\end{tabular}
\end{table}

For MWKR, remaining route work is estimated as
\[
W_{jo} = \sum_{q=o}^{last(j)} \min_{m \in A_{jq}} p_{jqm}.
\]

\end{document}
